% =========================================================================
\section{Attribution and algebraic approximation}\label{sec:attribution}

High local order and geometric iteration are not new principles. Traub organizes classical iteration
formulas and their error terms~\cite{traub61,traub64}. Gander discusses Halley's method and radical
corrections~\cite{gander}; Scavo and Thoo analyze its geometry~\cite{scavo}. Melman connects geometry
with global convergence, and Amat, Busquier and Guti\'errez develop geometric constructions for
nonlinear iteration~\cite{melman,amat}. Here AK is Newton under inversion and centered AD is Halley,
realized in the shared rectangle; their scalar formulas are classical. Reading all supports as
approximations of one function $\kappa_*$ (Lemma~\ref{lem:contact}) is the usual principle that the
order of an iteration is the order of contact of its iteration function with the exact one, specialized
to the rectangle.

The decentered and three-mean supports are algebraic corrections defined by explicit positive
radicals. Algebraic approximants generalize rational approximation and form an established
methodology~\cite{sergeev,shafer}, and the coefficient matching of Section~\ref{sec:mean} belongs to it;
Appendix~\ref{app:equiv} gives a specified quadratic comparison for both supports. The subdivision mean
proportional is Euclid's~\cite{euclid}; its final readout is an ordinary Pad\'e
approximant~\cite{dlmf,baker}, with the enclosure and width proved in Section~\ref{sec:enclosure}.
Counting particular constructions follows the concerns of g\'eom\'etrographie and
nomography~\cite{lemoine,ocagne}, and the efficiency index is Ostrowski's~\cite{ostrowski}.

The three-mean support is a linear combination of three members of the family
$q^{2}=t+a(t-1)^{2}$, and its method is not new either. Since $q_*$ is a homogeneous mean of $1$ and
$t$, the support belongs to a classical subject, the approximation and bounding of means such as the
logarithmic and extended means by combinations of the arithmetic, geometric and quadratic
means~\cite{carlson,stolarsky}. Its coefficient matching is of Pad\'e type in a wider sense: in the
variable $w$ it is a quadrature rule for the kernel $\sqrt{1+aw}$ with two prescribed nodes, in the way
ordinary Pad\'e approximants of Stieltjes series are Gaussian quadratures for the kernel
$1/(1+aw)$~\cite{baker}; and in the variable $\sqrt t+1/\sqrt t$ it is a quadratic algebraic approximant
subject to one constraint (Appendix~\ref{app:equiv}). Raising the order of a root iteration by an affine
combination of known iterations is also an established device, used for $n$th roots with rational
iterations in~\cite{dubeau}.

What this paper contributes is the stated rectangle realizations, their domains, the scalar convergence
proofs with explicit constants, the protocol bounds, and the classification of degrees.

% =========================================================================
\section{Evidence and reproducibility}\label{sec:evidence}

\begin{table}[htbp]
\centering\small
\caption{Status of each claim. ``Lean'': checked by the Lean kernel, in the environments described below.
``CAS'': exact identity between polynomials or rational functions, recomputed with a computer-algebra
system. ``Numerical'': high-precision evaluation, not a proof.}\label{tab:evidence}
\begin{tabularx}{\linewidth}{@{}l>{\raggedright\arraybackslash}p{0.19\linewidth}>{\raggedright\arraybackslash}X@{}}
\toprule
Claim & Domain & Status\\
\midrule
Lemma~\ref{lem:contact}, Table~\ref{tab:supports} & analytic supports & Lean (limit form); constants of AK, AD and projective AD: CAS\\
Theorem~\ref{thm:arc} & real $p>2$ & Lean (global convergence and the fifth-order limit); $O(e^{7})$ remainder written\\
Theorem~\ref{thm:mean} & real $p>2$ & Lean (positivity, global convergence, the limit $e_+/e^{9}\to c_M(p)$, positivity of $c_M$, $c_M(3)$); identities also CAS; $O(e^{11})$ remainder written\\
Remark~\ref{rem:family} & real $p>2$ & CAS (exact series)\\
Proposition~\ref{prop:rational} & rational $p$ & Lean, for both radical supports\\
Proposition~\ref{prop:scalar} & real $p\neq0$ (scalar) & written proof (change of variable in Theorem~\ref{thm:mean})\\
Proposition~\ref{prop:irrational} & real $p>0$ & Lean\\
Proposition~\ref{prop:mean-inventory} & integer $p\ge3$ & written inventory; total Lean; incidence simulation\\
Proposition~\ref{prop:index}, Table~\ref{tab:index} & real and integer $p\ge3$ & Lean; table values numerical\\
Tables~\ref{tab:cubic}, \ref{tab:p20} & stated examples & numerical\\
Proposition~\ref{prop:enclosure} & $X\ge1$, $p\ge1$ & written proof; scalar part in Lean (retained)\\
Proposition~\ref{prop:allp} & real $p\neq0$, $X>0$ & Lean (invariant, gap, inversion)\\
Appendix~\ref{app:equiv}, three-mean paragraph & $p=3$ & exact rational computation (series only)\\
\bottomrule
\end{tabularx}
\end{table}

\paragraph{Lean.}
The modules below are those of V28 and V29, unchanged in V30 and V31; they were compiled in the project
Docker environment (Lean~4.33.1, Mathlib~v4.33.1)~\cite{lean4,mathlib}, with no admitted proof, and their
main theorems depend only on \texttt{propext}, \texttt{Classical.choice} and \texttt{Quot.sound}. The
number of theorems of each module is given in parentheses.
\begin{itemize}
\item \texttt{DecenteredArcGlobal} (30): Theorem~\ref{thm:arc}, global monotone convergence of the
  decentered arc for every real $p>2$, by the sign argument of Section~\ref{sec:arc}.
\item \texttt{ContactLimit} (12): Lemma~\ref{lem:contact} in limit form, and the analytic tool used for
  the local orders: a finite sum $\sum_jc_je^{\lambda_jx}$ whose first $N$ moments
  $\sum_jc_j\lambda_j^{n}$ vanish is $O(x^{N})$.
\item \texttt{ArcOrderFive} (10): the fifth-order limit \eqref{eq:arc-order} with its constant.
\item \texttt{ThreeMeanGlobal} (57): parts (i) and (ii) of Theorem~\ref{thm:mean} for every real $p>2$,
  in the variable $\sqrt t$: the identities \eqref{eq:mean-pos}, \eqref{eq:mean-slope} and
  \eqref{eq:mean-square}, the positivity of the nine certificate polynomials, the enclosure of the
  correction between $1$ and $t^{1/p}$, global convergence, and the agreement of the Lean definition
  with \eqref{eq:qM}.
\item \texttt{ThreeMeanOrderNine} (29): part (iii), the limit $e_+/e^{9}\to c_M(p)$ for every real
  $p>2$, with $c_M>0$ and $c_M(3)=23/20736$. The radical is replaced by its quartic Taylor polynomial at
  an exact algebraic cost, and the remaining exponential sum has $22$ terms and ten vanishing moments
  (the tool of \texttt{ContactLimit}).
\item \texttt{ThreeMeanWork} (7): the total $10p+11$ of Table~\ref{tab:mean-inventory}; the inequality
  $\iota_M>\iota_{\mathrm{dec}}$ for every $p>0$, from $5^{11}<9^{9}$; and
  Proposition~\ref{prop:rational} for the three-mean support.
\item \texttt{DegreesEfficiency} (47), depending on Mathlib only: the exact contact identity;
  Propositions~\ref{prop:irrational} and~\ref{prop:index} for the classical and decentered supports;
  the trace totals of Section~\ref{sec:work}; the invariant of Proposition~\ref{prop:allp}.
\item \texttt{DegreesNested} (5): Proposition~\ref{prop:rational} for the decentered arc, with integer
  powers only.
\item \texttt{GeometricSubdivision}: scalar interval invariant, real-degree rational bracket and width
  of Section~\ref{sec:enclosure}; stored seven-declaration audit, not recompiled~\cite{supplS}.
\end{itemize}
\texttt{ArcOrderFive} and \texttt{ThreeMeanOrderNine} reuse the exponential-sum expansion of
$q_*^{2}/t$ and the width $w$ from the V29 modules \texttt{NestedOrderNine} and
\texttt{NestedRadicalReal}, and \texttt{DegreesNested} and \texttt{DegreesEfficiency} also contain the
nested-support statements of V29. These modules therefore remain in the build~\cite{supplNR}, but none of
their nested-support theorems is a result of V30 or V31. The local orders are stated as limits of
$e_+/e^{\,\omega+1}$; the refinements $O(e^{7})$ and $O(e^{11})$ of the remainders, which follow from
oddness, are written arguments.

\paragraph{Computer algebra and numerics.}
The accompanying script \texttt{verify\_v31.py} recomputes, with exact rational arithmetic: the
identities \eqref{eq:arc-der} (at sample points) and \eqref{eq:arc-id}; the constants of
Table~\ref{tab:supports}; the Pad\'e identity and the width equality \eqref{eq:width}; and the appendix
formulas. For the three-mean support it recomputes exactly, and independently of Lean: the coefficient
identities and the degenerate case $p=2$; \eqref{eq:mean-pos}, \eqref{eq:mean-slope} and
\eqref{eq:mean-square} with the positivity of the nine certificate polynomials; the series
\eqref{eq:mean-D4} with symbolic $h$ and the constant $c_M$; the series of Remark~\ref{rem:family}; and
the rational computation of Appendix~\ref{app:equiv}. Numerically it checks \eqref{eq:mean-slope}
against a numerical derivative, the limit $e_+/e^{9}\to c_M(p)$ at $p=2.5,3,4,7,20$, the enclosure
$\min(1,t^{1/p})<\Gamma<\max(1,t^{1/p})$ on $7200$ points with $p\in[2.001,10^{4}]$ and
$t\in[10^{-12},10^{12}]$, and the protocol of Section~\ref{sec:mean-construction}: the line and circle
incidences are evaluated at $60$ digits for $120$ configurations ($p$ from $3$ to $20$, four inputs,
five states), with the shape condition \eqref{eq:aspect} and, with one more join, without it, and
return the scalar update. It regenerates Tables~\ref{tab:cubic}, \ref{tab:index}, \ref{tab:p20} and
\ref{tab:depth} (working precision 300 to 2600 digits).

The earlier subdivision verifier checks three exact polynomial identities, 188 rational rail
incidences, four rational depth certificates, the outward-rounded $p=37.5$ example and 256 numerical
enclosures at 100 digits~\cite{supplS}. Finite numerical probes supplement proofs and do not certify
rounded implementations. The gallery contains companion interactive views; it does not certify finite
drawing error.

\paragraph{Figures.}
All figures are generated by \texttt{genfigs.py} as exact TikZ vector drawings from the formulas of the
paper; the script asserts, for every geometric drawing, that the constructed point $P_+$ agrees with the
scalar update to $10^{-12}$.

\paragraph{Revision history.}
The classical supports and the decentered arc were introduced in~\cite{v20} and restated with complete
proofs for a self-contained comparison. V28 added the contact lemma, the efficiency-index comparison and
the degree classification of Table~\ref{tab:degrees}; V29 added the three-mean support of
Section~\ref{sec:mean}; V30~\cite{v30} removed the nested-radical support of V27--V29 (retained in
V29~\cite{v29} and in the Lean repository), because the three-mean support attains the same order nine
with a smaller error constant, a protocol shorter by $48$ traces and three quartics instead of $345$
terms. The linear baseline, called MA in earlier versions, was renamed MB after its support line $MB$
(Appendix~\ref{app:mb}). The present revision (V31) preserves every mathematical result and formal proof
of V30 unchanged while reorganizing the exposition for readability: adding the opening visual overview
(Figure~\ref{fig:anatomy}), the target-contact plot (Figure~\ref{fig:target}), the step-by-step
three-mean storyboard (Figure~\ref{fig:mean}) and the efficiency charts (Figure~\ref{fig:efficiency}),
disambiguating local notation ($\lambda_X$ vs.\ $\lambda$), and moving the explicit degree-eight
certificate table into Appendix~\ref{app:certs}.

% =========================================================================
\section{Scope and limits}\label{sec:limits}

The following limits apply to the whole paper and are stated here once.
\begin{enumerate}[label=\arabic*.]
\item \textbf{Approximation.} Every calculator produces a convergent sequence or a finite enclosure. No
  root is claimed to be exactly constructible in finitely many operations.
\item \textbf{Finite steps.} The fast iterations are finite constructions for integer and rational
  degrees only, at a cost governed by the numerator $\mu$. For irrational $p$ the convergence theorems
  are scalar statements; the geometric calculator is the subdivision.
\item \textbf{Counts.} Trace counts are upper bounds for the specified protocols under the convention of
  Section~\ref{sec:convention}, which does not charge fixed preparation or decisions and weighs a join
  and an arc equally. They are not minimal, and they are not times, energies or drawing sizes. They are
  written inventories, not a formal cost semantics. The inventories do not use the convention to the
  same extent: the three-mean protocol takes its moving centres and compass openings from free
  intersections with seven prepared objects and prescribes the shape of the rectangle, whereas the
  decentered arc uses three prepared objects.
\item \textbf{Efficiency.} The index of Proposition~\ref{prop:index} is asymptotic and inherits the
  previous limit; a stricter convention on preparation would raise the cost of the three-mean
  support.
\item \textbf{Real drawings.} Exact incidences, the scalar real iteration and a physically inaccurate
  drawing are different objects. Conditioning, finite drawing error and floating-point behaviour are not
  analysed.
\item \textbf{Priority.} We have not identified the decentered or three-mean formulas with a prior formula,
  but a bounded comparison cannot establish priority, and different explicit formulas may be equivalent
  under changes of variable (Appendix~\ref{app:equiv}).
\item \textbf{Formalization.} Lean covers exactly the entries so marked in Table~\ref{tab:evidence}.
  The convergence theorems of the two radical supports and their local orders are formalized. Not
  formalized: the constants of the classical supports in Table~\ref{tab:supports}, the higher-order
  remainders, Remark~\ref{rem:family}, Proposition~\ref{prop:scalar} and Appendix~\ref{app:equiv}. For
  the inventories Lean checks the arithmetic of the totals, not the geometric protocols.
\end{enumerate}

% =========================================================================
\section{Conclusion and open questions}\label{sec:conclusion}

Changing one line in the Pandrosion rectangle moves the iteration from order two to order nine, and the
order is read off the contact of that line's slope with a single function $\kappa_*$. Two radical
supports extend the classical ones: a decentered arc of order five at the mobile cost of Halley's
construction, and a combination of three means of order exactly nine for every real $p>2$, globally
monotone, with constant $c_M(p)$ and a mobile cost of three arcs and one join. In trace counts, under
the stated conventions, the three-mean support has the largest efficiency index at every degree
$p\ge3$. These iterations are finite constructions for rational degrees; for a general real degree the
power-free subdivision with a Pad\'e readout gives a certified enclosure of relative width at most
$z^{5}/256$ in $15n+63$ traces.

Open questions:
\begin{enumerate}[label=\arabic*.]
\item Formalize a cost semantics for the drawing protocols, so that the inventories, and not only
  their totals, are machine-checked, and so that prepared objects can be weighed.
\item The family of Section~\ref{sec:mean} gives order seven with a single radical
  (Remark~\ref{rem:family}), but on a bounded range of residuals and with a state that crosses the
  root. Is there a single-radical support of order greater than five that is defined and monotone for
  every $t>0$? Which order do three conic means, or four, reach with positive radicands?
\item Analyse the two-level recursion in which the residual $u^{p}$ of an irrational degree is itself
  enclosed by subdivision: which inner accuracy preserves the outer order?
\item Analyse conditioning and finite drawing error.
\item Compare the exact formulas with earlier algebraic approximants and with the literature on means
  under specified transformations.
\item Extend to the three-mean support the half-plane contraction proved for the decentered arc in the
  complex companion~\cite{companion}. Numerically the boundary criterion of that paper holds for $p$
  from $2.0001$ to $10^{4}$, with $\sup_y\operatorname{Arg}\varphi(iy)\,p/\pi\approx0.503$ against the
  bound $1$; this is an observation, not a proof.
\end{enumerate}

% =========================================================================
\FloatBarrier
\appendix
\section{Classical rectangle supports}\label{app:classical}

\subsection{AK and centered AD: orders two and three}

\emph{AK: a fixed line.} Prepare $K=(W(1-X/p),0)$. The support $AK$ has slope $Hp/(WX)$, so
$\kappa=p$ and $\Tm_{\mathrm{AK}}(s)=ps/(p-1+t)$, Newton's method under inversion. It has order two and
converges from every positive start.

\emph{AD: a centered length transfer.} Prepare
\[ O_{\mathrm{AD}}=\Bigl(W-\frac{2W}{p-1},\;H-\frac{H(p+1)}{X(p-1)}\Bigr). \]
Transfer $AE=Hs^{p}$ downward from this centre to $D_{\mathrm{AD}}$ and join $AD_{\mathrm{AD}}$. Then
$\kappa=\tfrac12[p+1+(p-1)t]$, giving
\begin{equation}\label{eq:AD}
\Tm_{\mathrm{AD}}(s)=s\,\frac{p+1+(p-1)t}{p-1+(p+1)t}.
\end{equation}
This Halley map has cubic order and global monotone convergence. Earlier geometric interpretations
appear in~\cite{gander,scavo}; here the rail encoding and readout are shared with AK.

\begin{figure}[htbp]
\centering\figAK
\caption{Pandrosion AK, order two, at $p=3$, $X=2$, $s=3/4$ ($W=2$, $H=4$). The fixed support $AK$
(blue) replaces the MB support, while the triangle chain (grey) and the correction direction (green)
are retained. The inset magnifies the dashed window around $P$, $T$ and $P_+$.}\label{fig:ak}
\end{figure}

\begin{figure}[htbp]
\centering\figAD
\caption{Pandrosion AD, order three, at the same state. The orange length $AE$ is transferred onto the
prepared vertical through $O_{\mathrm{AD}}$; the lower arc point determines the support
$AD_{\mathrm{AD}}$ (blue). Here $O_{\mathrm{AD}}=C$.}\label{fig:ad}
\end{figure}

Both drawings retain the triangle stage, the green transport and the horizontal readout; only the
support changes.

\subsection{Projective AD: order four}

Two joins supply the support. Prepare the horizontal line and the point
\[ \ell:\ y=H\Bigl(1+\frac{5p-1}{(p+1)X}\Bigr),\qquad
Z_{\mathrm{proj}}=\Bigl(W\Bigl(1+\frac{5p-1}{2p(2p-1)}\Bigr),\;H\Bigl(1+\frac{p+1}{(2p-1)X}\Bigr)\Bigr). \]
Set $D_{\mathrm{proj}}=Z_{\mathrm{proj}}E\cap\ell$ and join $AD_{\mathrm{proj}}$. Substitution in
\eqref{eq:support} gives the $\kappa$ of Table~\ref{tab:supports} and
\begin{equation}\label{eq:proj}
\Tm_{\mathrm{proj}}(s)=s\,\frac{2p\bigl((2p-1)t+p+1\bigr)}{Q_{\mathrm{proj}}(t)},\qquad
Q_{\mathrm{proj}}(t)=(p+1)t^{2}+2(2p-1)(p+1)t+(2p-1)(p-1).
\end{equation}
The direct-root correction is the $[2/1]$ Pad\'e approximant; the reciprocal-state factor displayed
here is $[1/2]$. Direct differentiation gives
\[ \Tm_{\mathrm{proj}}'(s)=-\frac{2p(p^{2}-1)(2p-1)(t-1)^{3}}{Q_{\mathrm{proj}}(t)^{2}}, \]
so the map has its global maximum $\alpha$ at $s=\alpha$. Also
$\Tm_{\mathrm{proj}}(s)/s-1=(1-t)\bigl((p+1)t+5p-1\bigr)/Q_{\mathrm{proj}}(t)$. Thus a first step lands
at or below $\alpha$ and subsequent steps increase to $\alpha$; the order is exactly four. The
earlier preprint gives further dynamics and the original triangle-chain realization~\cite{v20}.

\begin{figure}[htbp]
\centering\figProjective
\caption{Projective AD, order four, at the same cubic state. The orange line $Z_{\mathrm{proj}}E$ meets
the prepared horizontal $\ell$ at $D_{\mathrm{proj}}$; the blue support $AD_{\mathrm{proj}}$, extended
below $A$, receives the green correction. The inset magnifies the dashed window and separates $P$, $T$
and $P_+$.}\label{fig:proj}
\end{figure}

\FloatBarrier
\section{The linear MB reference}\label{app:mb}

MB is a baseline, not an accelerated method. Its support is the line $MB$, from the prepared point $M$
to the lower-right vertex $B$; it does not pass through $A$. Earlier versions called this construction
MA: the name referred to the same lower-right vertex, which is $B$ in the labelling of this paper. The
parallel to $ME$ through $P$ meets $MB$, and a horizontal returns the next state. For $s\neq1$,
\[ \Tm_{\mathrm{MB}}(s)=1-\frac{X-1}{X\,\Sigma_p(s)},\qquad \Sigma_p(s)=\sum_{j=0}^{p-1}s^{j}. \]
For $X>1$ this map is increasing and
$\Tm_{\mathrm{MB}}(s)-s=(1-Xs^{p})/(X\Sigma_p(s))$, so it converges monotonically to $\alpha$ from every
positive start, with the removable extension at $s=1$. Its derivative at $\alpha$ is
$1-p\alpha^{p-1}/\Sigma_p(\alpha)\in(0,1)$: convergence is linear.

\begin{figure}[htbp]
\centering\figMB
\caption{Pandrosion MB, the linear reference construction, at $p=3$, $X=2$, $s=3/4$. Thin lines
construct the power by similar triangles. The green parallel meets the blue support $MB$; the
horizontal gives $P_+$. At $s=1$ the geometric intersection is coincident despite the removable scalar
extension.}\label{fig:mb}
\end{figure}

\FloatBarrier
\section{Equivalence questions for the radical corrections}\label{app:equiv}

Whether the arc is a known formula ``up to conjugacy'' depends on the class of coordinate changes
allowed. Let $\varphi$ denote the correction factor $\Phi_{\mathrm{arc}}(t)=(1+g)/(t+g)$ with
$g=(p-1)q_5$. Eliminating the radical gives the quadratic relation
\begin{equation}\label{eq:quad}
P_p(t,\varphi):=(p-1)^{2}S(t)(\varphi-1)^{2}-12p\,(t\varphi-1)^{2}=0 .
\end{equation}
The explicit radical in \eqref{eq:arc} selects the branch; the polynomial alone loses sign information
and has a double root at $t=1$.

\paragraph{Homographic changes.}
For each real $p>2$, $\Phi_{\mathrm{arc}}$ is not rational in $t$. Indeed $\disc_tS=96p(p+1)$ and
$S(0)=p-2>0$, so $S$ has two simple nonzero roots and is not a square in $\mathbb{C}(t)$; if the
correction were rational, so would be
$\sqrt{S(t)/(12p)}=(1-t\Phi_{\mathrm{arc}})/[(p-1)(\Phi_{\mathrm{arc}}-1)]$. Invertible homographies of
the input or output preserve rationality. For integer $p\ge3$, $S(y^{p})$ also has only simple roots,
so the state map $y\,\Phi_{\mathrm{arc}}(y^{p})$ is not rational and cannot be M\"obius-conjugate to a
rational state map, even as a germ at $y=1$. This state-map exclusion is asserted only for integer
degrees.

\paragraph{A specified quadratic normalization.}
A diagonal type $(2,2,2)$ quadratic Hermite--Pad\'e (Shafer) form for $f(t)=t^{-1/p}$ requires
polynomials of degree at most two with $A(t)f^{2}+B(t)f+C(t)=O((t-1)^{8})$~\cite[eqs.~(1.4)--(1.6)]{cabay}.
Direct substitution into \eqref{eq:quad} gives instead, with $\xi=t-1$,
\[ P_p\bigl(1+\xi,(1+\xi)^{-1/p}\bigr)=\frac{(p-2)(p-1)^{2}(p+1)(3p+1)}{60p^{4}}\,\xi^{6}+O(\xi^{7}), \]
contact six, not eight. The coefficients of $P_p$ have no common factor and its discriminant in
$\varphi$ is $48p(p-1)^{2}(t-1)^{2}S(t)$, so it is irreducible over $\mathbb{C}(t)$; any polynomial
with the same degree bounds defining this branch is a constant multiple of $P_p$. This excludes that
particular diagonal approximant in these coordinates, without excluding other algebraic families or
transformed targets.

\paragraph{Local analytic conjugacy.}
Distinct formulas may share a local analytic normal form: B\"ottcher's theorem conjugates a
superattracting germ $F(e)=ce^{\varrho}+O(e^{\varrho+1})$, $c\neq0$, to its monomial normal form~\cite{buff}. This
supplies neither a finite incidence protocol nor a global domain or drawing inventory. The explicit
construction and its proved dynamics are therefore distinct from a claim of a new local conjugacy class.

\paragraph{The three-mean correction as a constrained quadratic approximant.}
Put $\sigma=\sqrt t+1/\sqrt t$, so that $w=\sigma^{2}-4$. Then
\[ \frac{q_5}{\sqrt t}=\sqrt{1+a_5(\sigma^{2}-4)},\qquad
\frac{q_M}{\sqrt t}=m\,\sigma+g-n\sqrt{1+d\,(\sigma^{2}-4)} , \]
and $y=q_M/\sqrt t$ satisfies $(y-m\sigma-g)^{2}=n^{2}\bigl(1+d(\sigma^{2}-4)\bigr)$, a quadratic
relation whose coefficients have degrees $0,1,2$ in $\sigma$. In this variable the three-mean support
is therefore a quadratic algebraic approximant of $q_*/\sqrt t$ at $\sigma=2$ in the sense
of~\cite{shafer}, but not the Hermite--Pad\'e approximant of its type: its discriminant is constrained to
have no term of degree one in $\sigma$. The constraint is what keeps the radicand a quadratic
polynomial in $t$, that is, one arc. The unconstrained form has one more free coefficient. For $p=3$
it is
\[ \frac{16}{33}-\frac{\sigma}{66}+\frac1{99}\sqrt{5\sigma^{2}+736\sigma+1424}
=\frac{q_*}{\sqrt t}+\frac{(\sigma-2)^{5}}{19131876}+O\bigl((\sigma-2)^{6}\bigr), \]
one order more than \eqref{eq:mean-D4}, where the error is $\tfrac{23}{354294}(\sigma-2)^{4}$; but
multiplied by $t$ its radicand contains $\sqrt t$, so the corresponding support is a nested radical.
This is an exact rational computation on the series at $p=3$; the corresponding iteration has not been
analysed.

% =========================================================================
\FloatBarrier
\section{Algebraic certificates for Theorem~\ref{thm:mean}}\label{app:certs}

Table~\ref{tab:mean-cert} lists the explicit polynomials $\Psi_1$, $\Psi_2$ and $\Psi_3$ appearing in
\eqref{eq:mean-slope} and \eqref{eq:mean-square} in the proof of Theorem~\ref{thm:mean}. Each coefficient
of $\Psi_2$ and $\Psi_3$ is checked both in Lean (\texttt{ThreeMeanGlobal}) and in \texttt{verify\_v31.py}
by expanding in powers of $p-2$, where all nine terms have strictly positive coefficients.

\begin{table}[htbp]
\centering\small
\caption{The certificate polynomials of Theorem~\ref{thm:mean}. $\Psi_1$ is given by its coefficients on
$\tau^{6}+1$, $\tau^{5}+\tau$, $\tau^{4}+\tau^{2}$, $\tau^{3}$; $\Psi_2$ and $\Psi_3$ by their
coefficients on $(\tau-1)^{4}$, $\tau(\tau-1)^{2}$, $\tau^{2}$. Each coefficient of $\Psi_2$ and
$\Psi_3$ is positive for $p>2$: written in powers of $p-2$ it has nine positive
coefficients.}\label{tab:mean-cert}
\begin{tabular}{@{}lll@{}}
\toprule
$\Psi_1$ & $\tau^{6}+1$ & $1175p^{8}+2755p^{6}+2352p^{4}+880p^{2}+128$\\
 & $\tau^{5}+\tau$ & $5120p^{8}-3584p^{6}-3264p^{4}+1984p^{2}-256$\\
 & $\tau^{4}+\tau^{2}$ & $-6635p^{8}+13877p^{6}+17808p^{4}+1808p^{2}-128$\\
 & $\tau^{3}$ & $75776p^{8}-121856p^{6}+54912p^{4}-9344p^{2}+512$\\
\midrule
$\Psi_2$ & $(\tau-1)^{4}$ & $30185p^{8}+32563p^{6}+16572p^{4}+8032p^{2}+128$\\
 & $\tau(\tau-1)^{2}$ & $54180p^{8}+100044p^{6}+85680p^{4}+32256p^{2}$\\
 & $\tau^{2}$ & $262836p^{8}-335160p^{6}+310464p^{4}$\\
\midrule
$\Psi_3$ & $(\tau-1)^{4}$ & $489845p^{8}+647956p^{6}+361344p^{4}+123904p^{2}-1024$\\
 & $\tau(\tau-1)^{2}$ & $804380p^{8}+1583824p^{6}+2050176p^{4}+1033216p^{2}-4096$\\
 & $\tau^{2}$ & $3829896p^{8}-3081456p^{6}+2225664p^{4}+2128896p^{2}$\\
\bottomrule
\end{tabular}
\end{table}

% =========================================================================
\begin{thebibliography}{99}\small\setlength{\itemsep}{1.5pt plus 0.5pt}
\bibitem{pappus} Pappus of Alexandria. \emph{Book III of the Mathematical Collection}. English
translation by John B. Little, following Hultsch; College of the Holy Cross, Bookshelf 63, 2023.
\url{https://crossworks.holycross.edu/hc_books/63/}.
\bibitem{launay} Micka\"el Launay. Pandrosion et la duplication du cube. Micmaths, video presentation,
2024. \url{https://www.youtube.com/watch?v=7T4neM660SM}. Rectangle construction and correction,
approximately 16:37--19:10.
\bibitem{sergeev} Alexei V. Sergeev and David Z. Goodson. Summation of asymptotic expansions of
multiple-valued functions using algebraic approximants: application to anharmonic oscillators.
\emph{J. Phys. A: Math. Gen.} 31:4301--4317, 1998.
\url{https://avsergeev1.narod.ru/papers_37_iop_paper.pdf}.
\bibitem{euclid} Euclid. \emph{Elements}, Book VI, Proposition 13. Online edition by David E. Joyce,
Clark University. \url{https://mathcs.clarku.edu/~djoyce/elements/bookVI/propVI13.html}. Accessed
17 October 2026.
\bibitem{gallery} Ivan Besevic. Real exponents: geometric subdivision. Pandrosion Geometry Workshop,
companion interactive construction, 2026.
\url{https://ivan-fr.github.io/pandrosion/gallery/geometric-subdivision.html}. Local source snapshot
included with the V26-S supplement.
\bibitem{dlmf} NIST Digital Library of Mathematical Functions. Section 3.11(iv): Pad\'e
approximations. \url{https://dlmf.nist.gov/3.11#iv}. Accessed 17 October 2026.
\bibitem{baker} George A. Baker, Jr. and Peter Graves-Morris. \emph{Pad\'e Approximants}, 2nd ed.
Encyclopedia of Mathematics and its Applications 59, Cambridge University Press, 1996.
doi:10.1017/CBO9780511530074. In particular Section~1.5, Theorems 1.5.2--1.5.3, on argument and
value homographic invariance.
\bibitem{traub61} J. F. Traub. On a class of iteration formulas and some historical notes.
\emph{Comm. ACM} 4(6):276--278, 1961. doi:10.1145/366573.366606.
\bibitem{traub64} J. F. Traub. \emph{Iterative Methods for the Solution of Equations}. Prentice-Hall,
1964.
\bibitem{ostrowski} A. M. Ostrowski. \emph{Solution of Equations and Systems of Equations}. Academic
Press, New York, 1960.
\bibitem{gander} Walter Gander. On Halley's iteration method. \emph{Amer. Math. Monthly}
92(2):131--134, 1985. doi:10.2307/2322644.
\bibitem{scavo} T. R. Scavo and J. B. Thoo. On the geometry of Halley's method. \emph{Amer. Math.
Monthly} 102(5):417--426, 1995. doi:10.2307/2975033.
\bibitem{melman} A. Melman. Geometry and convergence of Euler's and Halley's methods. \emph{SIAM
Review} 39(4):728--735, 1997. doi:10.1137/S0036144595301140.
\bibitem{amat} S. Amat, S. Busquier and J. M. Guti\'errez. Geometric constructions of iterative
functions to solve nonlinear equations. \emph{J. Comput. Appl. Math.} 157(1):197--205, 2003.
doi:10.1016/S0377-0427(03)00420-5.
\bibitem{shafer} R. E. Shafer. On quadratic approximation. \emph{SIAM J. Numer. Anal.}
11(2):447--460, 1974. doi:10.1137/0711037.
\bibitem{lemoine} \'Emile Lemoine. \emph{G\'eom\'etrographie ou Art des constructions
g\'eom\'etriques}. Scientia 18, C. Naud, Paris, 1902.
\url{https://catalogue.bnf.fr/ark:/12148/cb36049032t}.
\bibitem{ocagne} Maurice d'Ocagne. \emph{Esquisse d'ensemble de la nomographie}. M\'emorial des
sciences math\'ematiques 4, Gauthier-Villars, 1925.
\url{https://www.numdam.org/item/MSM_1925__4__1_0/}.
\bibitem{lean4} Leonardo de Moura and Sebastian Ullrich. The Lean 4 theorem prover and programming
language. In \emph{CADE 28}, LNCS 12699, pages 625--635. Springer, 2021.
doi:10.1007/978-3-030-79876-5\_37.
\bibitem{mathlib} The mathlib Community. The Lean mathematical library. In \emph{CPP 2020}, pages
367--381. ACM, 2020. doi:10.1145/3372885.3373824.
\url{https://leanprover-community.github.io/papers/mathlib-paper.pdf}.
\bibitem{supplNR} Ivan Besevic. Nested-radical Pandrosion support: derivation, geometric example and
Lean convergence certificates. Source bundle accompanying \emph{Geometric Root Calculators}, V27-NR,
18 October 2026. Includes \texttt{NestedRadicalCubic}, \texttt{NestedRadicalRealCertificate},
\texttt{NestedRadicalReal}, \texttt{NestedRadicalAllDegrees}, audit scripts, source hashes and
numerical checks.
\bibitem{supplS} Ivan Besevic. Geometric subdivision and Pad\'e readout: supplementary source and
verification records. Source bundle accompanying \emph{Geometric Root Calculators}, V26-S, 2026.
Includes \texttt{GeometricSubdivision.lean}, its retained axiom audit and the independent subdivision
verifier.
\bibitem{v30} Ivan Besevic. Geometric root calculators, version 30. Research preprint, 16 November
2026. \url{https://github.com/ivan-fr/pandrosion/blob/main/paper/Geometric_Root_Calculators_v30.pdf}.
\bibitem{v29} Ivan Besevic. Geometric root calculators, version 29. Research preprint, 13 November
2026. \url{https://github.com/ivan-fr/pandrosion/blob/main/paper/Geometric_Root_Calculators_v29.pdf}.
\bibitem{v20} Ivan Besevic. Reciprocal root geometry, version 20. Research preprint, 2026.
\url{https://github.com/ivan-fr/pandrosion/tree/main/paper/reciprocal_geometry_v20}.
\bibitem{cabay} Stan Cabay and George Labahn. A fast, reliable algorithm for calculating
Pad\'e--Hermite forms. Technical Report CS-89-09, University of Waterloo, 1989.
\url{https://cs.uwaterloo.ca/research/tr/1989/CS-89-09.pdf}. Equations (1.4)--(1.6), printed page 1,
PDF page 5.
\bibitem{buff} Xavier Buff, Adam L. Epstein and Sarah Koch. B\"ottcher coordinates, 2011.
\url{https://arxiv.org/abs/1104.2981}. Introduction: the classical one-variable theorem.
\bibitem{carlson} B. C. Carlson. The logarithmic mean. \emph{Amer. Math. Monthly} 79:615--618, 1972.
\bibitem{stolarsky} Kenneth B. Stolarsky. Generalizations of the logarithmic mean. \emph{Math. Mag.}
48(2):87--92, 1975. doi:10.2307/2689825.
\bibitem{dubeau} Fran\c{c}ois Dubeau. Newton's method and high-order algorithms for the $n$th root
computation. \emph{J. Comput. Appl. Math.} 224:66--76, 2009. doi:10.1016/j.cam.2008.04.014.
\bibitem{companion} Ivan Besevic. Complex geometric root calculators, version C3. Research preprint,
companion to this paper, 13 November 2026.
\end{thebibliography}
